\section{Entorno de trabajo en Python sobre MSTM}
\label{sec:python_wrapper}

El código MSTM~\cite{Mackowski_code} se controla mediante un fichero
de texto plano (\texttt{mstm.inp}) con sintaxis numérica de Fortran,
devuelve los resultados en ficheros \texttt{.dat} cuya estructura
depende de las opciones activadas, y trabaja exclusivamente en unidades
adimensionales escaladas por el número de onda. Preparar a mano cada
fichero de entrada, convertir las unidades, lanzar el ejecutable y
procesar las salidas es tedioso y propenso a errores, sobre todo cuando
se quieren repetir muchas simulaciones parecidas. Para automatizar todo
ese ciclo se ha desarrollado un entorno de trabajo en Python con
interfaz gráfica.\footnote{El código completo está disponible en
\url{https://github.com/GorkaLarrazabalEhu/mie-scattering-project}.}

El entorno se diseñó para poder explorar distintas configuraciones y
observables, no solo los que se emplean en este trabajo. Por eso
incluye más opciones de las que luego aparecen en los resultados; en
esta sección se describe el conjunto general y, cuando procede, se
indica qué parte se usa en cada caso. El desarrollo de esta herramienta
ha constituido una parte sustancial del trabajo.

El entorno se organiza en una interfaz gráfica, un módulo de utilidades
físicas y un conjunto de scripts de representación:

\begin{itemize}
    \item \texttt{THE\_MAN.py}: interfaz gráfica (GUI) basada en
          Tkinter que centraliza la entrada de parámetros, gestiona el
          proyecto y orquesta la ejecución y el postproceso.
    \item \texttt{mstm\_utils.py}: utilidades físicas; conversión a
          unidades adimensionales e interpolación del índice de
          refracción a la longitud de onda deseada.
    \item Scripts de representación (\texttt{plot\_nearfield.py},
          \texttt{plot\_scattering\_matrix.py},
          \texttt{plot\_radius\_sweep.py} y
          \texttt{plot\_asymmetry.py}): generan las figuras de campo
          cercano, campo lejano, barridos y parámetro de asimetría.
\end{itemize}

% ---------------------------------------------------------------------
\begin{figure}[H]
    \centering
    \includegraphics[width=1\textwidth]{Imagenes/Pantallazo_GUI.png}
    \caption{Interfaz gráfica del \emph{wrapper} en Python para MSTM,
    configurada para una simulación de ejemplo: cuatro esferas de oro de
    radio 0.1\,$\mu$m dispuestas en un cuadrado a 45$^\circ$ respecto a la
    onda incidente, con un barrido en longitud de onda de 400 a 800\,nm en
    30 pasos. Se aprecia también la vista previa del agregado, que se
    actualiza en tiempo real.}
    \label{GUI:screenshot}
\end{figure}

% ---------------------------------------------------------------------
\subsection{Arquitectura y flujo de trabajo}
\label{subsec:arquitectura}

Toda la configuración de una simulación se almacena en un único objeto
(clase \texttt{Config}) que agrupa la geometría del agregado, el haz
incidente, el material, las opciones de campo cercano y los barridos.
Ese objeto se serializa en un fichero JSON, de modo que cualquier
simulación queda descrita por completo en un fichero de texto legible.

El ciclo de trabajo habitual, que se sigue íntegramente desde la
interfaz de la figura~\ref{GUI:screenshot}, consta de cuatro pasos:

\begin{enumerate}
    \item El usuario introduce los parámetros en la GUI.
    \item Al pulsar \texttt{Write mstm.inp}, la GUI convierte las
          unidades físicas a las adimensionales de MSTM, da formato de
          Fortran a los números y escribe el fichero de entrada con la
          sintaxis adecuada.
    \item \texttt{Run MSTM Simulation} lanza el ejecutable
          \texttt{mstm.exe} mediante el módulo \texttt{subprocess} de
          Python.
    \item Los botones \texttt{Plot} leen los ficheros \texttt{.dat}
          resultantes y generan las figuras.
\end{enumerate}

La comunicación entre Python y Fortran se realiza estrictamente por
archivos: Python escribe la entrada y lee las salidas, sin acoplarse al
código de Fortran. Como el binario puede tardar en converger, su
ejecución se lanza en un hilo aparte y se muestra una ventana de
progreso mientras dura; de este modo el programa no se cuelga ni deja de
responder. Cuando la simulación termina, se informa del resultado. 

Cada simulación se guarda en su propia carpeta, que contiene el fichero
de entrada, las salidas y la configuración en JSON. Esto mantiene los
resultados autocontenidos y reproducibles: una simulación guardada puede
recargarse más tarde y la GUI reconstruye a partir del JSON la geometría
y todos los parámetros, lista para volver a ejecutarla o para representar
sus resultados. Todos los resultados de este trabajo se han generado
mediante esta interfaz.

% ---------------------------------------------------------------------
\subsection{Parámetros físicos y unidades adimensionales}
\label{subsec:mstm_utils}

MSTM trabaja con todas las longitudes escaladas por el número de onda en
el medio externo, $k_{0} = \frac{2\pi}{\lambda}$,
de modo que cada esfera queda caracterizada por su parámetro de tamaño
$x = k_{0}\,a$, siendo $a$ el radio, siguiendo la convención de la
sección~\ref{sec:miecoeffs}. El módulo \texttt{mstm\_utils.py} convierte
los parámetros geométricos del usuario (longitud de onda y radio en
$\mu$m) al par (\texttt{length\_scale\_factor}, \texttt{size\_parameter})
que requiere el fichero de entrada.

El usuario puede introducir directamente un índice de refracción
complejo o seleccionar un material de una pequeña base de datos. En este
último caso, el módulo carga tablas de
\texttt{refractiveindex.info} almacenadas localmente (Au, Ag, Cu y Fe a
partir de las medidas de Johnson \& Christy; Al según Rakić; agua según
Hale \& Querry) e interpola linealmente el índice complejo
$\tilde{n}$ a la longitud de onda solicitada. 
% El espaciado
% de estas tablas ($\sim 10$\,nm) es suficientemente fino para que el error
% de interpolación sea despreciable frente a la incertidumbre experimental
% de los datos. 
Cabe señalar que MSTM adopta el convenio de parte
imaginaria positiva para medios absorbentes, convención que el módulo
respeta al leer las tablas.

% ---------------------------------------------------------------------
\subsection{Definición de la geometría}
\label{subsec:geometria}

La geometría del agregado se especifica en una tabla de texto en la que
cada fila describe una esfera mediante su posición $(x,y,z)$, su radio y
su índice de refracción. De este modo se pueden construir tanto agregados
homogéneos como mezclas de materiales o tamaños, sin más que ajustar los
valores de cada fila.

Para los casos más habituales, la GUI puede generar automáticamente las
posiciones de varias geometrías sencillas: dos o cuatro esferas en línea
(en dirección norte--sur o este--oeste), cuatro esferas en cuadrado y
cuatro en cuadrado girado $45^\circ$ respecto a la onda incidente. Las
posiciones se calculan a partir del radio y de la separación entre
superficies $s$ que indica el usuario: la distancia centro a centro de un
par de esferas se fija como $d = a_i + a_j + s$, de manera que el hueco
entre superficies es exactamente $s$ con independencia de los radios. Se
puede asignar un radio común a todas las esferas o un radio distinto a
cada una.

La GUI muestra además una vista previa del agregado en el plano $XZ$
(figura~\ref{GUI:screenshot}) que se actualiza en tiempo real a medida
que se modifican las posiciones, los radios o la separación, lo que
facilita comprobar la disposición antes de lanzar la simulación.

% ---------------------------------------------------------------------
\subsection{Barridos paramétricos}
\label{subsec:sweeps}

El entorno permite repetir una misma configuración variando un
parámetro, ya sea la longitud de onda o el radio de las esferas. En lugar
de invocar el binario una vez por cada valor, la GUI genera un único
\texttt{mstm.inp} con bloques \texttt{new\_run} sucesivos, aprovechando el
mecanismo nativo de MSTM. Los dos barridos son excluyentes: solo uno puede
estar activo a la vez.

En un barrido en longitud de onda, cada \texttt{new\_run} actualiza el
factor de escala $k_{0} = 2\pi/\lambda$ y el índice de refracción
$\tilde{n}(\lambda)$ interpolado, manteniendo fija la geometría física del
agregado. El resultado son las eficiencias y demás observables en función
de $\lambda$.

En un barrido en radio, la longitud de onda es fija, de modo que $k_{0}$
y $\tilde{n}$ son constantes y se calculan una sola vez; lo único que
cambia entre \emph{runs} es el radio de las esferas. Cuando se usa una de
las geometrías predefinidas, los centros se recalculan en cada paso con la
misma regla $d = 2r + s$, de manera que la separación entre superficies
$s$ se mantiene constante mientras crece el radio. Esto permite estudiar
cómo cambia la respuesta del agregado con el tamaño de las partículas sin
alterar el hueco entre ellas.

% ---------------------------------------------------------------------
\subsection{Muestreo del campo cercano}
\label{subsec:adaptive_step}

El paso espacial $\Delta$ de la malla del campo cercano debe resolver
simultáneamente la longitud de onda, el radio de la esfera más pequeña y
el \emph{gap} entre esferas adyacentes. El módulo calcula automáticamente
un paso candidato como el mínimo de tres criterios:

\begin{equation}
    \Delta = \min\!\left(
        \frac{2\pi}{N_{\text{ppw}}},\;
        \frac{r_{\min}}{N_{\text{rad}}},\;
        \frac{g_{\min}}{N_{\text{gap}}}
    \right),
    \label{eq:dmin}
\end{equation}

Cada criterio fija un número mínimo de puntos a lo largo de una escala
característica del problema. $N_{\text{ppw}}$ (\emph{points per
wavelength}) es el número de puntos por longitud de onda: como en
unidades escaladas una longitud de onda equivale a $2\pi$, el término
$2\pi/N_{\text{ppw}}$ garantiza que las oscilaciones del campo queden bien
muestreadas. $N_{\text{rad}}$ es el número de puntos a lo largo del radio
de la esfera más pequeña, $r_{\min}$, y $N_{\text{gap}}$ el número de
puntos a lo largo del hueco más estrecho entre dos esferas.
% $g_{\min} = \min_{i\neq j}\!\bigl(\|\mathbf{r}^i - \mathbf{r}^j\| - a^i -
% a^j\bigr)$. 
Este último criterio es importante porque es justo en esos
huecos donde el campo varía de forma más brusca, y donde aparecen los
fenomenos que interesa analizar. Tomar el mínimo de los tres
asegura que se muestrea adecuadamente la más fina de esas escalas. Los
valores por defecto son $N_{\text{ppw}}=50$, $N_{\text{rad}}=25$ y
$N_{\text{gap}}=12$; además, el paso se acota a un rango razonable y, si
la malla resultante superase $2\times 10^{6}$ puntos, se reescala hasta
satisfacer esa cota.

El cálculo del campo cercano es costoso y puede tardar bastante, sobre
todo con mallas finas, por lo que conviene no elegir un paso demasiado
pequeño a ciegas. Este criterio no pretende dar el paso definitivo, sino
una estimación del orden de magnitud al que conviene trabajar: en este
trabajo se ha usado como punto de partida y, tras inspeccionar los
resultados, se ha empleado ese mismo valor o uno ligeramente menor. El
usuario puede en cualquier momento sustituir el valor automático por uno
manual.

% ---------------------------------------------------------------------
\subsection{Visualización de resultados}
\label{subsec:visualizacion}

Cada tipo de figura se genera con un script independiente que la GUI lanza
como un proceso aparte, pasándole el fichero de salida correspondiente.
Esto mantiene la representación separada del resto del programa y permite
abrir varias figuras a la vez.

\subsubsection*{Campo cercano}

El script \texttt{plot\_nearfield.py} lee el fichero de campo cercano de
MSTM, que contiene en cada punto de la malla las componentes complejas de
los campos eléctrico y magnético para las dos polarizaciones de la onda
incidente, tal como se describe en la sección~\ref{sec:interpretacion}. En
cada punto se dispone así de 24 componentes: para cada una de las dos
polarizaciones (paralela y perpendicular), los campos eléctrico y
magnético, cada uno con sus tres componentes espaciales ($x,y,z$) y cada
componente con su parte real e imaginaria ($2\times2\times3\times2=24$); el
script puede representar cualquiera de ellas. A partir de los campos
calcula además la intensidad $|\mathbf{E}|^2$ y, para luz no polarizada,
su promedio incoherente entre polarizaciones
(ec.~\ref{eq:nearfield_intensity}), así como el vector de Poynting
promedio. Los resultados se representan como mapas de contorno sobre la
malla reconstruida, con las secciones transversales de las esferas
superpuestas.

\subsubsection*{Campo lejano}

El script \texttt{plot\_scattering\_matrix.py} extrae del fichero de
salida la matriz de Mueller y las eficiencias de extinción, absorción y
dispersión de cada \emph{run}. Representa la función de fase
$S_{11}(\theta)$ en coordenadas polares (en escalas lineal y logarítmica
de forma simultánea) y las eficiencias $Q_{\text{ext}}$, $Q_{\text{abs}}$
y $Q_{\text{sca}}$ frente a la longitud de onda cuando hay barrido
espectral, o frente al radio cuando el barrido es en tamaño
(\texttt{plot\_radius\_sweep.py}).

El programa calcula también los grados de polarización del campo
dispersado a partir de sus parámetros de Stokes
(ec.~\ref{eq:MuellerMatrix}, \ref{eq:stokes_farfield}) y el parámetro de 
asimetría $g$
(ec.~\ref{eq:g_mie}, mediante \texttt{plot\_asymmetry.py}). Son magnitudes
disponibles en el entorno que se han empleado de forma puntual según el
caso.